\documentclass[11pt]{article}

\usepackage[margin=25mm]{geometry}

\usepackage{amsmath,amssymb,fontspec,unicode-math}

\IfFontExistsTF{DejaVu Serif}{\setmainfont{DejaVu Serif}}{\setmainfont{Latin Modern Roman}}

\IfFontExistsTF{Latin Modern Math}{\setmathfont{Latin Modern Math}}{\setmathfont{DejaVu Math TeX Gyre}}

\usepackage{xurl,hyperref}

\setlength{\parindent}{0pt}\setlength{\parskip}{6pt}

\begin{document}

\title{Every Rollout Counts: Optimal Resource Allocation for Efficient Test-Time Scaling}

\author{ELVIA preliminary AI analysis}\date{21 October 2026}\maketitle

Preliminary assessment: Major.\par

Independent scientific review remains pending. Author adoption has not been requested.

Exact Beta and embedding examples defeat the theoretical bridge to optimal success allocation and the claimed exact direction recovery. Independent rounding and released pre-pruning add finite-budget discrepancies; empirical search results are not remeasured.

\section*{Review scope}

Rapid claim-focused reading of current author-paper passages and relevant released source. Direct mathematical or static-source reasoning is stated below. No author program, inherited verifier, model inference or benchmark was executed.

\section*{Problem 1: A log-utility calculation does not prove the advertised success optimum}

Source: Proceedings Eq.(7), PDF p.4; Proposition 2, p.6; Theorem 1, p.7; Appendix B.2, p.17

Assessment: supported; impact: major

\subsection*{Source claim (paraphrase)}

The paper’s independent Beta model minimizes failure \(F(B)=\prod_i\prod_{r=0}^{B_i-1}[\kappa(1-w_i)+r]/(\kappa+r)\), but its proportional-allocation justification switches to weighted logarithmic utility.

\subsection*{Analysis}

For \(k=2,N=3,\kappa=10,w=(2/3,1/3)\), proportional allocation \((2,1)\) has success \(271/297\), while \((3,0)\) has \(839/891\), an improvement of \(26/891\). This disproves unrestricted finite-concentration optimality, separately from the stated low-confidence regime.

In that regime, take \(k=N=2,w=(9/10,1/10),\kappa=1/100\). Nearest-rounded proportional counts \((2,0)\) give success \(9099/10100\); coverage \((1,1)\) gives \(91/100\), greater by \(92/10100\). As \(\kappa\downarrow0\), these tend to \(0.90\) and \(0.91\). Proper Beta laws and allowed quality scores realize both examples.

The limit rewards covering distinct candidates; if \(N\geq k\), any integer allocation covering all candidates can be optimal. By contrast, \(\sum_jQ_j\log B_j\) genuinely has proportional optimizer, and its KL gap is correct for that utility. No identity turns that gap into a success-probability comparison.

\subsection*{Finding and scope}

The central Bayesian-optimality bridge and success interpretation are unsupported; the log-utility result and empirical DORA comparisons remain distinct.

\subsection*{Proposed correction}

The existing exact integer optimizer repeatedly chooses a coordinate maximizing \(w_i/(\kappa+B_i)\). Each next rollout reduces log failure by \(-\log[1-\kappa w_i/(\kappa+B_i)]\); these marginal gains decrease along each coordinate, so greedy selection of the largest remaining prefix gain is optimal. This is a result for the independent Beta model, not evidence that actual LLM rollouts obey it.

The existing correction is retained as a bounded proposal; author adoption is not established.

\section*{Problem 2: Finite-temperature affinity does not exactly cancel direction multiplicity}

Source: Proceedings Theorem 1, PDF p.7; Eqs.(10)–(12), pp.6–7; Appendix B.3, pp.17–18

Assessment: supported; impact: major

\subsection*{Source claim (paraphrase)}

Theorem 1 claims exact recovery of direction-level allocation from the affinity diagonal when scores are equal within each direction, without a separation assumption.

\subsection*{Analysis}

Use unit embeddings \(e_1=e_2=(1,0)\), \(e_3=(0.99,\sqrt{1-0.99^2})\), zero quality scores and \(T_s=0.01\). Write \(a=e\). Row-softmax diagonals are \(a/(2a+1),a/(2a+1),a/(a+2)\). The normalized total masses are \[W_A=\frac{2(a+2)}{4a+5},\qquad W_B=\frac{2a+1}{4a+5}.\] Equal direction scores require \((1/2,1/2)\), but \(W_A-1/2=3/[2(4a+5)]>0\). The contradiction occurs before rounding.

Actual inverse cluster sizes arise only under extra conditions, such as identical within-cluster embeddings and a zero-temperature limit with strict cross-cluster separation. Finite positive cross-affinities prevent exact inverse-size cancellation in general.

\subsection*{Finding and scope}

The exact recovery theorem fails under its written assumptions; the heuristic may still improve real search.

\subsection*{Proposed correction}

The supplied restricted result assumes uniqueness weights \(\gamma_i=c/k_j\) within direction \(j\), with a common \(c>0\). Summing quality–uniqueness products then cancels \(k_j\) exactly. A soft approximation requires a quantitative separation/error bound.

The existing correction is retained as a bounded proposal; author adoption is not established.

\section*{Problem 3: Independent rounding and quality pre-pruning change the finite allocation}

Source: Proceedings Eq.(13), PDF p.7; current author scripts/sal/search/dora.py, allocation lines 290–310 and next-iteration beam construction

Assessment: supported; impact: minor

\subsection*{Source claim (paraphrase)}

The fixed rollout constraint is \(\sum_iB_i=N\), but the displayed rule rounds each normalized allocation independently. The release also removes zero quality-rounded candidates before diversity weighting.

\subsection*{Analysis}

Weights \((0.4,0.4,0.2)\) and \(N=4\) give rounded counts \((2,2,1)\), totaling \(5\). They are realized by orthogonal embeddings and quality scores \((\log2,\log2,0)\). Later truncating a repeated beam list restores its length but makes the allocation depend on order.

With three identical direction-A embeddings and one orthogonal direction-B embedding, quality weights \((0.3,0.3,0.3,0.1)\) round at budget \(4\) to \((1,1,1,0)\), discarding B before uniqueness. Applying quality and inverse-size diversity together instead yields roughly equal final weights and \((1,1,1,1)\). The current source has the quality-only pruning stage; no full author search was run.

\subsection*{Finding and scope}

The displayed/released finite allocation differs from the constrained rule; effects on published model-search accuracy are unmeasured.

\subsection*{Proposed correction}

For a proportional heuristic, compute all intended weights first and use floor counts plus largest fractional remainders under a declared tie rule to preserve the budget. That restores feasibility, not optimal success; the Beta optimizer is a separate correction.

The existing correction is retained as a bounded proposal; author adoption is not established.

\section*{Limits}

Preliminary AI source/argument assessment; independent human scientific adjudication remains pending.\par

Exact calculations and static-source findings do not remeasure published experiments or establish the historical experiment-generating revision.\par

\section*{Source identity}

Current NeurIPS 2025 proceedings PDF; freshly retrieved; latest applicable listed author revisions and resource heads compared

\url{https://proceedings.neurips.cc/paper_files/paper/2025/file/940b66bcdb34fb9ae47d85de838c372d-Paper-Conference.pdf}

Reviewed PDF SHA-256: \nolinkurl{8469eb452c639efb258d238163521a9132efa8d97830c750f4f783cd2b0988da}

The exact inspected PDF was freshly retrieved from the official author/publication endpoint. Its SHA-256 identifies inspected bytes; PDF creation metadata is not treated as publication chronology. Authors and title were checked against the primary header. Historical raw reports remain unchanged outside this export.

\end{document}
